% Options for packages loaded elsewhere
\PassOptionsToPackage{unicode}{hyperref}
\PassOptionsToPackage{hyphens}{url}
%
\documentclass[
  twoside,
  12pt,
  letterpaper,
  fleqn]{article}

\usepackage{amsmath,amssymb}
\usepackage{lmodern}
\usepackage{iftex}
\ifPDFTeX
  \usepackage[T1]{fontenc}
  \usepackage[utf8]{inputenc}
  \usepackage{textcomp} % provide euro and other symbols
\else % if luatex or xetex
  \usepackage{unicode-math}
  \defaultfontfeatures{Scale=MatchLowercase}
  \defaultfontfeatures[\rmfamily]{Ligatures=TeX,Scale=1}
\fi
% Use upquote if available, for straight quotes in verbatim environments
\IfFileExists{upquote.sty}{\usepackage{upquote}}{}
\IfFileExists{microtype.sty}{% use microtype if available
  \usepackage[]{microtype}
  \UseMicrotypeSet[protrusion]{basicmath} % disable protrusion for tt fonts
}{}
\makeatletter
\@ifundefined{KOMAClassName}{% if non-KOMA class
  \IfFileExists{parskip.sty}{%
    \usepackage{parskip}
  }{% else
    \setlength{\parindent}{0pt}
    \setlength{\parskip}{6pt plus 2pt minus 1pt}}
}{% if KOMA class
  \KOMAoptions{parskip=half}}
\makeatother
\usepackage{xcolor}
\setlength{\emergencystretch}{3em} % prevent overfull lines
\setcounter{secnumdepth}{-\maxdimen} % remove section numbering
% Make \paragraph and \subparagraph free-standing
\ifx\paragraph\undefined\else
  \let\oldparagraph\paragraph
  \renewcommand{\paragraph}[1]{\oldparagraph{#1}\mbox{}}
\fi
\ifx\subparagraph\undefined\else
  \let\oldsubparagraph\subparagraph
  \renewcommand{\subparagraph}[1]{\oldsubparagraph{#1}\mbox{}}
\fi


\providecommand{\tightlist}{%
  \setlength{\itemsep}{0pt}\setlength{\parskip}{0pt}}\usepackage{longtable,booktabs,array}
\usepackage{calc} % for calculating minipage widths
% Correct order of tables after \paragraph or \subparagraph
\usepackage{etoolbox}
\makeatletter
\patchcmd\longtable{\par}{\if@noskipsec\mbox{}\fi\par}{}{}
\makeatother
% Allow footnotes in longtable head/foot
\IfFileExists{footnotehyper.sty}{\usepackage{footnotehyper}}{\usepackage{footnote}}
\makesavenoteenv{longtable}
\usepackage{graphicx}
\makeatletter
\def\maxwidth{\ifdim\Gin@nat@width>\linewidth\linewidth\else\Gin@nat@width\fi}
\def\maxheight{\ifdim\Gin@nat@height>\textheight\textheight\else\Gin@nat@height\fi}
\makeatother
% Scale images if necessary, so that they will not overflow the page
% margins by default, and it is still possible to overwrite the defaults
% using explicit options in \includegraphics[width, height, ...]{}
\setkeys{Gin}{width=\maxwidth,height=\maxheight,keepaspectratio}
% Set default figure placement to htbp
\makeatletter
\def\fps@figure{htbp}
\makeatother


\usepackage{geometry}

\geometry{reset, letterpaper, height=9in, width=6in, hmarginratio=1:1, vmarginratio=1:1, marginparsep=0pt, marginparwidth=0pt, headheight=15pt}

\usepackage{lipsum}

\def\dklogo{\raisebox{-.2\height}{\includegraphics[width=12pt]{../../assets/direct-knowledge-logo-book.png}}\footnotesize{\ Direct Knowledge}}
\usepackage{fancyhdr}
\pagestyle{fancy}
\fancyhead[LE]{\thepage}%
\fancyhead[RE]{\dklogo}%
\fancyhead[CE]{}%
\fancyfoot[LE]{}%
\fancyfoot[RE]{}%
\fancyfoot[CE]{}%
\fancyhead[LO]{D. A. Smith}%
\fancyhead[RO]{\thepage}%
\fancyhead[CO]{}%
\fancyfoot[LO]{}%
\fancyfoot[RO]{}%
\fancyfoot[CO]{}%


\renewcommand{\headrulewidth}{0.4pt}
\renewcommand{\footrulewidth}{0pt}

\usepackage{setspace}
\linespread{1.25}

\usepackage{amssymb}
\usepackage{amsmath}
\usepackage{amsthm}
\usepackage{graphicx} 


\usepackage[most]{tcolorbox} 
\definecolor{block-gray}{gray}{0.97}
\newtcolorbox{zitat}[1][]{%
    colback=block-gray,
    grow to right by=-10mm,
    grow to left by=-10mm, 
    boxrule=0pt,
    boxsep=0pt,
    breakable,
    enhanced jigsaw,
    borderline west={2pt}{0pt}{gray},
    colbacktitle={block-gray},
    coltitle={black},
    fonttitle={\large\bfseries},
    attach title to upper={},
    #1,
}
\renewcommand{\quote}{\zitat}
\renewcommand{\endquote}{\endzitat}

\def\banner{ {\rule{\linewidth}{0.2pt}}
    \begin{minipage}[c]{14px}\includegraphics[width=14pt]{../../assets/direct-knowledge-logo-book.png}\end{minipage}
    \begin{minipage}[c]{150px}Direct Knowledge | Blog \end{minipage}
}

\usepackage{titling}
\setlength{\droptitle}{-8ex}
\pretitle{\vspace{-20pt}\banner\vspace{10pt}\begin{flushleft}\huge\bfseries}
\posttitle{\par\end{flushleft}}
\preauthor{\begin{flushleft}\Large}
\postauthor{,
    \small{Direct Knowledge, USA}
    \footnote{email: david@directknowledge.com}
    \footnote{\copyright \, 2023 \ David A. Smith}
    \footnote{With authorization, you can freely share and reproduce portions of this work for educational or personal use. Please note that distributing any portion of it in print form requires further permission from its original authors, as does posting online to public servers or mailing lists without prior consent.}
    \end{flushleft}
    }
\predate{\begin{flushleft}}
\postdate{\end{flushleft}}

\renewenvironment{abstract}
{\par\noindent\textbf{\abstractname.}\ \ignorespaces \itshape}
{\par\medskip}


\makeatletter
\makeatother
\makeatletter
\makeatother
\makeatletter
\@ifpackageloaded{caption}{}{\usepackage{caption}}
\AtBeginDocument{%
\ifdefined\contentsname
  \renewcommand*\contentsname{Table of contents}
\else
  \newcommand\contentsname{Table of contents}
\fi
\ifdefined\listfigurename
  \renewcommand*\listfigurename{List of Figures}
\else
  \newcommand\listfigurename{List of Figures}
\fi
\ifdefined\listtablename
  \renewcommand*\listtablename{List of Tables}
\else
  \newcommand\listtablename{List of Tables}
\fi
\ifdefined\figurename
  \renewcommand*\figurename{Figure}
\else
  \newcommand\figurename{Figure}
\fi
\ifdefined\tablename
  \renewcommand*\tablename{Table}
\else
  \newcommand\tablename{Table}
\fi
}
\@ifpackageloaded{float}{}{\usepackage{float}}
\floatstyle{ruled}
\@ifundefined{c@chapter}{\newfloat{codelisting}{h}{lop}}{\newfloat{codelisting}{h}{lop}[chapter]}
\floatname{codelisting}{Listing}
\newcommand*\listoflistings{\listof{codelisting}{List of Listings}}
\makeatother
\makeatletter
\@ifpackageloaded{caption}{}{\usepackage{caption}}
\@ifpackageloaded{subcaption}{}{\usepackage{subcaption}}
\makeatother
\makeatletter
\@ifpackageloaded{tcolorbox}{}{\usepackage[many]{tcolorbox}}
\makeatother
\makeatletter
\@ifundefined{shadecolor}{\definecolor{shadecolor}{rgb}{.97, .97, .97}}
\makeatother
\makeatletter
\makeatother
\ifLuaTeX
  \usepackage{selnolig}  % disable illegal ligatures
\fi
\usepackage[citestyle = authoryear]{biblatex}
\addbibresource{../references.bib}
\IfFileExists{bookmark.sty}{\usepackage{bookmark}}{\usepackage{hyperref}}
\IfFileExists{xurl.sty}{\usepackage{xurl}}{} % add URL line breaks if available
\urlstyle{same} % disable monospaced font for URLs
\hypersetup{
  pdftitle={Writing Mathematics},
  pdfauthor={David A. Smith},
  pdfkeywords={mathematics, writing, publishing},
  hidelinks,
  pdfcreator={LaTeX via pandoc}}

\title{Writing Mathematics}
\usepackage{etoolbox}
\makeatletter
\providecommand{\subtitle}[1]{% add subtitle to \maketitle
  \apptocmd{\@title}{\par {\large #1 \par}}{}{}
}
\makeatother
\subtitle{How I Write Mathematics}
\author{David A. Smith}
\date{Saturday, January 28, 2023}

\begin{document}
\maketitle
\begin{abstract}
Writing mathematics clearly, carefully, and accurately is a meaningful
skill for people of all ages to learn and understand. It requires
considerable focus, such as understanding concepts with specialized
language and notation. Writing mathematics can be difficult, but it is
worth pursuing because it helps people comprehend complex ideas and
prepares them for advancing their own.
\end{abstract}
\ifdefined\Shaded\renewenvironment{Shaded}{\begin{tcolorbox}[borderline west={3pt}{0pt}{shadecolor}, enhanced, interior hidden, boxrule=0pt, sharp corners, breakable, frame hidden]}{\end{tcolorbox}}\fi

\hypertarget{introduction}{%
\subsection{Introduction}\label{introduction}}

One of the greatest challenges for students is learning to express
mathematics fluently and accurately. Trying to decipher complex formulas
or explaining mathematical ideas in writing can be a daunting task, even
for those with an advanced understanding of mathematics. However,
mastering this skill can result in tremendous rewards. Expressing
mathematical ideas effectively through written language provides us with
a greater understanding and appreciation of the subject, not only from
an academic perspective but from its relevance to our daily lives as
well. It's important to understand that although making sense of
mathematics on paper may seem difficult at first, it's worth pursuing
because developing this proficiency can help tremendously in class
discussions and exams while also providing insights into real-world
problems we might face on a regular basis.

\hypertarget{what-makes-writing-mathematics-different}{%
\subsection{What Makes Writing Mathematics
Different?}\label{what-makes-writing-mathematics-different}}

Writing mathematics is an artful combination of two distinct languages,
natural and mathematics. Natural language offers the writer numerous
expression possibilities but can be ambiguous. In contrast, writing with
mathematical symbols requires clarity yet allows succinctness when
conveying complex concepts. Accurate language is a necessity when it
comes to mathematics. To properly communicate your thoughts and
opinions, ensure that each word carries the correct connotation you
intend for them to have.

\begin{quote}
The precision of words is essential in math-based conversations.
\end{quote}

To comprehend mathematically expressed ideas, readers must take their
time and read content several times while allowing themselves ample
pauses for contemplation along the way. Further, mathematics writing
often serves as reference material, which necessitates that its contents
should ideally be accessed piece by piece upon demand rather than
require deep immersion into text volumes at once.

Writing mathematics has rules -- some narrow, others broad. Small
conventions relate to sentence structure (including punctuation) and are
easily verifiable. Broader ones involve general style and strategies
that depend on the author's discretion. For a full primer on
mathematical writing see \textcite{krantz2017primer}.

Crafting an effective document takes more than just following a few
simple rules; it requires mastering composition. This skillfulness is
where the art of writing comes in, with broad and deep strategies for
your entire work. It's all about connection: linking individual
sentences together for clarity and flow while adhering to specific
sentence structure regulations like commas or mathematical
terminologies. At its most complex level, these interrelated
requirements can create intricate webs that are both precise yet
engaging -- and (if followed) can result in a masterpiece.

Now it's time to dive into each of \textbf{the five levels for writing
mathematics}: sentence, paragraph, section, document, and beyond.

\hypertarget{sentence-level}{%
\subsection{Sentence Level}\label{sentence-level}}

Breaking complex mathematical ideas into more accessible parts is
essential for effective communication. Using the right voice and symbols
in your writing is crucial for success.

\begin{quote}
These small details can make comprehension natural -- but mastering them
is fundamental when creating potent written works.
\end{quote}

Here are ten rules that I use when writing mathematics.

\begin{enumerate}
\def\labelenumi{\arabic{enumi}.}
\tightlist
\item
  Do not start a sentence with a mathematical symbol.
\item
  Do not end a sentence with a mathematical symbol, instead use the
  correct punctuation.
\item
  Replace logical symbols with words, unless you are writting on logic.
\item
  Do not use a colon at the end of a sentence. Do not use sentence
  fragments, just use complete statements only.
\item
  Do spell out small numbers, when they are not being used
  mathematically.
\item
  Use ``we'' (``you'' and ``me'' together) for formal exposition, use
  ``I'' (``you'' and ``I'') for informal discussions.
\item
  Use ``that'' to indicate a specific object (a restrictive, essential
  clause) and not for making a point of inference.
\item
  Use ``which'' to add information to objects (a nonrestrictive,
  nonessential clause) and not as a conjunction.
\item
  Separate formulas with words. Do not list formulas, but rather
  communicate with a reader.
\item
  Write a sentence that flows logically from left to right to eliminate
  confusion.
\end{enumerate}

There are two obvious observations to make about these rules. First, it
is impossible to delineate a complete list of rules for natural
language. So think of these rules given here not as a complete list, but
only as a guide that you should (strongly) consider following when
writing mathematics.

The second observation is that these rules do not take into account
mathematical accuracy. Of course, your mathematics should be accurate.
This is not up for debate because mathematics that is not accurate is
simply not mathematics. Even a slight exchange of words can lead to
impotence or even nonsense. For example, consider the following
variation of mathematical induction.

\begin{quote}
If \(A\) is a subset of the natural numbers \(\mathbb{N}\) with the
following two properties: \(i\)) 0 is in \(A\) and \(ii\)) for all \(k\)
in \(\mathbb{N}\), \(k\in A\) implies \(k+1\in A\), then \(A\) is the
set of natural numbers.
\end{quote}

Now let's switch two words and see what you think.

\begin{quote}
If \(A\) is \textbf{the} subset of the natural numbers \(\mathbb{N}\)
with the following two properties: \(i\)) 0 is in \(A\) and \(ii\)) for
all \(k\) in \(\mathbb{N}\), \(k\in A\) implies \(k+1\in A\), then \(A\)
is \textbf{a} set of natural numbers.
\end{quote}

In mathematics, precision is everything. For more on these rules see
\textcite{knuth1989mathematical}.

\hypertarget{paragraph-level}{%
\subsection{Paragraph Level}\label{paragraph-level}}

Crafting the perfect opening paragraph is essential to grabbing a
reader's attention and inviting them on an engaging journey through your
prose. To do so, divide your work into linear segments to be presented
in a hierarchical development for ease of comprehension. Retain an even
flow by keeping notation and terms familiar while avoiding lengthy
explanations or eloquent phrasing that could confuse readers.
Furthermore, provide previews of upcoming topics before diving deeper
into the content so they can prepare themselves mentally as you write.
Adding visuals where appropriate will also aid understanding -- use
examples and counterexamples along with suggestive references when
necessary to maximize clarity without sacrificing creativity.

\begin{enumerate}
\def\labelenumi{\arabic{enumi}.}
\tightlist
\item
  Begin a paragraph with your best sentence to ensure that readers stay
  hooked.
\item
  Each sentence needs its own distinct beat. Peruse your writing and
  fine-tune it until everything flows smoothly.
\item
  A theorem (definition, lemma, corollary, etc.) should stand firmly on
  its own two feet and not rely upon what came before.
\item
  Pay particular attention to terms like ``therefore'', as they are
  essential for achieving the right cadence and has many variations such
  as ``whence'',``hence'', ``and so'', etc.
\item
  Break up paragraphs by displaying important formulas on a line by
  themselves.
\item
  Don't over complicate things with unnecessary subscripts,
  superscripts, or other vertically spaced symbols.
\item
  Break up consecutive long sentences into simpler ones, and break up
  mathematics (but not important formulas) into readable text.
\item
  Don't be tempted to resort to the use of technical jargon without good
  reason.
\item
  Every paragraph should have a (mathematical) point, make it clear.
\item
  Write a paragraph that flows logically from sentence to sentence to
  eliminate confusion.
\end{enumerate}

Mathematics writing requires you to effectively communicate your thought
process and convince the reader that your conclusions are valid. A
first-rate mathematical exposition should provide precise explanations
while also being persuasive enough to satisfy a skeptical audience.

\hypertarget{section-level}{%
\subsection{Section Level}\label{section-level}}

Crafting a clear and concise section requires consistent information,
deliberate structure, and reader-minded readability.

\begin{enumerate}
\def\labelenumi{\arabic{enumi}.}
\item
  \textbf{Commmunicate}. While writing, the writer must envision what
  areas of confusion could arise for their reader and strive to ensure
  that these are addressed. This foresight is crucial in order to
  effectively communicate with readers; otherwise it risks creating
  misunderstandings instead of a meaningful exchange. To ensure accuracy
  and clarity, having an audience in mind while writing is not simply
  beneficial but necessary.
\item
  \textbf{Examples}. Mastering the art of example and counterexample can
  be a powerful tool. However, it is important to not overuse this power
  -- examples should have some ``spark'' or element that will engage
  readers with thought-provoking insight. Additionally, don't forget the
  importance of providing context by using both examples and
  counterexamples which illustrate definitions/results while
  simultaneously clarifying any underlying essential assumptions.
\item
  \textbf{Figures}. Whenever possible, facilitate understanding by
  utilizing visual elements that have the power to communicate complex
  ideas quickly and accurately. Design figures with captions as
  reinforcement for concepts discussed in textual form; this will help
  bring key points into clear focus without being bogged down in
  minutiae. Work diligently to utilize graphs over tables when
  illustrating main topics.
\item
  \textbf{Efficiency}. When writing proofs, strive for efficiency by
  taking advantage of earlier results and avoiding unnecessary
  repetition. If a current proof seems to be significantly similar to
  one already constructed in the past, think beyond simple alterations;
  try finding an interesting overarching generalization that ties both
  together.
\item
  \textbf{Strategize}. When crafting especially intricate arguments,
  provide strategic comments throughout your construction -- noting the
  purpose behind each step taken as well as providing overviews can make
  lengthy demonstrations easier to follow. When tackling a complex
  proof, take the time to strategize before you write. Start by
  outlining how you plan to approach it and give periodic updates on
  your progress. At the appropiate time and place, demonstrate why
  particular steps are necessary and look for similarities between
  results in order to avoid reiterating similar arguments.
\end{enumerate}

\begin{quote}
Be strategically efficient.
\end{quote}

Crafting a perfect section for your chapter is no easy feat. It requires
the careful placement of each element -- from preliminary discussions to
lemmas, proofs and more -- into an order that both makes sense logically
while also being effectively digested by readers. The ideal outcome
presents detail in such a way so as to capture attention, yet surprise
them with unexpected connections between elements. All parts fit
together like pieces of a puzzle, forming an explanation or argument far
greater than its individual components.

\begin{enumerate}
\def\labelenumi{\arabic{enumi}.}
\tightlist
\item
  While writing, the writer must envision what areas of confusion could
  arise for their reader and strive to ensure that these are addressed.
  This foresight is crucial in order to effectively communicate with
  readers; otherwise it risks creating misunderstandings instead of a
  meaningful exchange. To ensure accuracy and clarity, having an
  audience in mind while writing is not simply beneficial but necessary.
\item
  Mastering the art of example and counterexample can be a powerful
  tool. However, it is important to not overuse this power -- examples
  should have some ``spark'' or element that will engage readers with
  thought-provoking insight. Additionally, don't forget the importance
  of providing context by using both examples and counterexamples which
  illustrate definitions/results while simultaneously clarifying any
  underlying essential assumptions.
\item
  Whenever possible, facilitate understanding by utilizing visual
  elements that have the power to communicate complex ideas quickly and
  accurately. Design figures with captions as reinforcement for concepts
  discussed in textual form; this will help bring key points into clear
  focus without being bogged down in minutiae. Work diligently to
  utilize graphs over tables when illustrating main topics.
\item
  When writing proofs, strive for efficiency by taking advantage of
  earlier results and avoiding unnecessary repetition. If a current
  proof seems to be significantly similar to one already constructed in
  the past, think beyond simple alterations; try finding an interesting
  overarching generalization that ties both together.
\item
  When crafting especially intricate arguments, provide strategic
  comments throughout your construction -- noting the purpose behind
  each step taken as well as providing overviews can make lengthy
  demonstrations easier to follow. When tackling a complex proof, take
  the time to strategize before you write. Start by outlining how you
  plan to approach it and give periodic updates on your progress. At the
  appropiate time and place, demonstrate why particular steps are
  necessary and look for similarities between results in order to avoid
  reiterating similar arguments.
\end{enumerate}

\hypertarget{document-level}{%
\subsection{Document Level}\label{document-level}}

Establishing an introduction that introduces the important theorems,
provides insight into what matters, and previews what's to come can
``hook'' readers in for further exploration. Make sure your document has
these components.

\textbf{Research}. To effectively communicate your mathematics, it's
essential to take into account the background knowledge of your intended
audience. Do the research necessary to know your audience. This may mean
teaching classes, going to workshops, collaborating on documents,
refereeing journals, or even reading education journals that specialize
in learning about your audience. While you are accomplishing this
research, keep in mind the following three points.

Firstly, don't let unfamiliar terminology confuse them; provide
additional explanation as needed for clarity and comprehension.
Secondly, consider adding appendices when discussing more complex or
specialized topics that may require deeper context in order to be fully
understood by less experienced readers. Thirdly, presenting too much
detail can be overwhelming for those without an expert's understanding;
however, avoiding explanations altogether risks leaving readers feeling
lost or confused.

\begin{quote}
Speak to someone and know that someone as well as you can.
\end{quote}

\textbf{Organize}. A well-crafted document requires a sound strategy. To
achieve this, start by outlining the expected sections in your paper or
chapter so you know exactly what to include in each step of its
development. Establish necessary definitions that will inform and
structure every result and proof before testing them for accuracy with
lemmas. Next, dabbling in some examples along the way allows readers to
apply concepts practically. Repeat these steps until there is an obvious
logical flow from beginning to end, then add interstitial comments as
needed. To a reader, these comments are often the linchpin needed. They
should provide further context on certain arguments, compare opposing
points within the literature or simply warn against any potential
missteps.

Make sure to keep your reader informed throughout; start each segment
with a brief introduction and perhaps an outline of what's ahead. Avoid
stretches without any clear direction. Instead, provide hints as you
progress by stating outcomes such as ``It follows that \(\ldots\)''.
Finally, be sure to reiterate essential points in conclusion for easy
recall later on.

When composing equations and propositions, don't make it difficult for
readers to find the references; clearly number each item along with
providing a content/name descriptor. To ensure concentration isn't
broken by excessive page flipping, repeat any simple math expressions
used in addition to reminding the reader of earlier analysis or unusual
notations they may have encountered. Don't be afraid to reiterate
information if necessary -- however remain mindful that too much
repetition can become exhausting.

\begin{quote}
Say something in the right place, at the right time, in the right
manner.
\end{quote}

Of course, the researching of your audience and the organizing of your
document go hand-in-hand. Or as I like to say:

\begin{quote}
Research your audience and organize for them.
\end{quote}

\hypertarget{and-beyond}{%
\subsection{And Beyond}\label{and-beyond}}

By reflecting my own personal tastes, this article offers unique
insights into the art of mathematical writing. Although some may not
agree or share similar opinions, I believe many people will find it
valuable and meaningful. While it may not satisfy all readers looking to
learn more about this topic, I trust that many have come here seeking a
connection between their passions and mine. If you are looking for more
on this topic see \textcite{halmos1970write},
\textcite{steenrod1973write}, and \textcite{higham1998handbook}.

\begin{quote}
While it might be hard to get started, writing mathematics is an
invaluable skill that requires patience and dedication. \textbf{But if
you put in the work your mastery of mathematical exposition can be a
force of change that improves lives}.
\end{quote}

Finally, I'd like to mention that this is where many beginners believe
math ends -- with books and articles (documents). But writing your
research program, (turning your suspicions, conjectures, unproven ideas,
and theories into written words) should be started as soon as possible.
Turning your ideas into written mathematics can enable you to harness
this force of change.


\printbibliography



\thispagestyle{empty}


\end{document}
